\documentclass[11pt]{article}
\usepackage[a4paper,margin=1in]{geometry}
\usepackage[T1]{fontenc}
\usepackage{lmodern,microtype}
\usepackage{amsmath,amssymb,amsthm,mathtools}
\usepackage{booktabs,enumitem}
\usepackage{xcolor}
\usepackage[numbers,sort&compress]{natbib}
\usepackage[colorlinks=true,linkcolor=blue!55!black,citecolor=blue!55!black]{hyperref}

\newtheorem{theorem}{Theorem}[section]
\newtheorem{proposition}[theorem]{Proposition}
\newtheorem{corollary}[theorem]{Corollary}
\theoremstyle{remark}
\newtheorem{remark}[theorem]{Remark}

\title{Expanded-Window Anchored Zonotope Obstruction}
\author{Bin Wang}
\date{July 2026}

\begin{document}
\maketitle

\begin{abstract}
RH-254 doubles the resolved candidate margin and supplies 16 new roots at
each archived endpoint.  We now perform the required anchored reachability
audit rather than a zero-target fit.  After discarding incomplete boundary
pairs, every conjugate shell receives a single-use weight $0\le w_j\le1$.
The resulting convex zonotope has zero passes at all 32 endpoints.  Its
weighted distance from the deterministic target ranges from $0.143585$ to
$0.423998$, with primal--dual gaps below $5.83\times10^{-15}$.  Hence every
prefix and each of the 62,030,604,700 eligible binary subsets also fails.
This is a finite margin-32 obstruction, not a global nonexistence theorem.
\end{abstract}

\section{Anchored expanded class}

For orders $n=2,\dots,12$, let $a=(a_n)$ be the deterministic target of
RH-243 \citep{WangRH243}.  At a fixed noise/channel endpoint, remove the
Perron and parity traces from the full matrix powers and write the remaining
base vector as $b$.  If $S_j$ is a conjugate-complete shell from RH-254, put
\begin{equation}
 v_j=\left(\sum_{\mu\in S_j}\mu^n\right)_{n=2}^{12}.
 \label{eq:shell}
\end{equation}
All $v_j$ are real up to the archived conjugacy tolerance.

The expanded single-use class is the zonotope
\begin{equation}
 \mathcal Z=\left\{b-\sum_jw_jv_j:0\le w_j\le1\right\}.
 \label{eq:box}
\end{equation}
Its anchored distance is measured by
\begin{equation}
 d(\mathcal Z,a)=\min_{0\le w\le1}
 \sum_{n=2}^{12}\frac{|b_n-a_n-(Vw)_n|}{n}.
 \label{eq:distance}
\end{equation}

\begin{proposition}[Box dominance]
\label{prop:dominance}
Every shell prefix, contiguous shell interval, and binary single-use subset is
contained in the box class \eqref{eq:box}.  Therefore
$d(\mathcal Z,a)>\varepsilon$ excludes all such selectors at tolerance
$\varepsilon$.
\end{proposition}

\begin{proof}
Each listed discrete selector has a weight vector with coordinates in
$\{0,1\}$, hence lies in $[0,1]^J$.  Minimizing over the larger box can only
decrease the distance.
\end{proof}

\section{Primal--dual audit}

The weighted $\ell^1$ problem \eqref{eq:distance} is a linear program.  We
solve both its box primal and explicit zonotope dual, then compare their
objectives.  The shell-complete ranks range from 33 to 64; no incomplete
conjugate fragment enters the optimization \citep{WangRH254}.

\begin{center}
\small
\begin{tabular}{lr}
\toprule
quantity & value\\
\midrule
endpoints & 32\\
eligible prefixes & 836\\
eligible binary subsets & 62,030,604,700\\
prefix passes & 0\\
box passes & 0\\
box distance range & 0.143585--0.423998\\
distance/tolerance range & 10.17--117.66\\
minimum failure margin & 0.142335\\
maximum primal--dual gap & $5.83\times10^{-15}$\\
\bottomrule
\end{tabular}
\end{center}

All 32 expanded distances improve relative to the frozen RH-248 box
\citep{WangRH248}; the improvement ranges from
$2.43\times10^{-7}$ to $0.00583514$.  Improvement is not reachability: the
best expanded point remains more than ten local tolerances away.

\begin{corollary}[Scoped expanded-window obstruction]
At each archived endpoint, no prefix or binary subset of the margin-32
conjugate-complete shells reaches the deterministic anchor within the local
noise tolerance.
\end{corollary}

\begin{proof}
The computed box distance is larger than tolerance at all endpoints, and
Proposition~\ref{prop:dominance} applies.  The dual objectives agree with the
primal values to the stated floating tolerance.
\end{proof}

\section{Boundary and next route}

The conclusion is confined to the finite margin-32 single-use shell class.
It does not exclude a still larger resolved window, unbounded weights, or a
signed/complex selector forced by invariant quotient algebra.  It also does
not provide interval enclosures for the eigensolver roots.

The live next route is therefore structural rather than another box scan:
derive signed or complex coefficients from a quotient identity and audit
whether they form a legal selector.  Gates A--E remain false/open.  No
Hilbert--Polya operator, Riemann-zero identification, zeta-divisor equality,
or RH implication is claimed.

\bibliographystyle{plainnat}
\bibliography{references}
\end{document}
